\chapter{Number theory and combinatorics}\label{ch:ntcomb}
\index{gcd}\index{sieve of Eratosthenes}\index{modular arithmetic}\index{binary exponentiation}\index{combinatorics}

Contests count on these two toolkits constantly: number theory for
divisibility, primes, and modular machinery; combinatorics for the
counting models every probability, DP, or hashing problem eventually
reduces to. After this chapter you can factorize at sieve speed,
move safely inside a modulus, choose between stars-and-bars and
inclusion--exclusion by the shape of the constraint, and spot the
linear recurrences that exponentiate.

\section{Divisibility and gcd}\label{sec:gcd}\index{Euclidean algorithm}
The Euclidean algorithm is an invariant machine: $\gcd(a,b)$ equals
$\gcd(b, a \bmod b)$ because common divisors of the pair are
unchanged by the reduction. The \emph{extended} version returns
B\'ezout coefficients --- $x,y$ with $ax+by=\gcd(a,b)$ --- which is
what makes modular division well-defined.
\lstinputlisting{code/gcdext.h}

\section{Primes}\label{sec:primes}\index{smallest prime factor}
The sieve marks every composite by its smallest factor; storing that
factor per cell gives factorization in $O(\log n)$ per query.
Trial division to $\sqrt{n}$ suffices for single values.
\begin{example}[CSES \emph{Counting Divisors}]\label{ex:cdiv}
Count the divisors of each queried $x$.
\par\smallskip
\textbf{Strategy.} SPF table once, then
$\prod(e_i+1)$ per factorization --- no $\sqrt{x}$ loop per query.
\par\smallskip
\begin{solution}
\lstinputlisting{code/counting_divisors.cpp}
Real output (CSES sample):
\begin{lstlisting}[numbers=none]
5
2
6
\end{lstlisting}
\end{solution}
\end{example}
\begin{figure}[ht]\centering
\begin{tikzpicture}[font=\tiny]
\foreach \i in {2,...,11}{\node[draw=slate!50,minimum size=0.32cm] at (\i*0.36,0) {\i};}
\node[font=\tiny,anchor=west] at (0.6,-0.6) {2 marks evens};
\node[font=\tiny,anchor=west] at (0.6,-1) {3 marks 6,9,12\dots};
\node[font=\tiny,anchor=west] at (0.6,-1.4) {survivors are prime};
\end{tikzpicture}

\caption{Sieve marking on $n=30$. Notice each prime marks only its
own multiples --- the survivors were never touched.}
\label{fig:sieve}
\end{figure}

\section{Modular arithmetic}\label{sec:mod}\index{modular inverse}
Addition, subtraction, multiplication and (with an inverse)
division all commute with reduction --- so reduce at every step.
The inverse exists iff $\gcd(a,m)=1$: for prime $m$,
$a^{-1}=a^{m-2}$ (Fermat); generally, extended gcd returns it.
\begin{example}[CSES \emph{Exponentiation II}]\label{ex:exp2}
$a^{b^c} \bmod 10^9{+}7$ for triples.
\par\smallskip
\textbf{Strategy.} The exponent itself is huge; Fermat reduces it
mod $m{-}1$ first (when $\gcd(a,m)=1$), then ordinary binpow.
\par\smallskip
\begin{solution}
\lstinputlisting{code/exponentiation2.cpp}
Real output (CSES sample):
\begin{lstlisting}[numbers=none]
2187
279632277
91039580
\end{lstlisting}
\end{solution}
\end{example}

\section{Binary exponentiation}\label{sec:binpow}\index{binpow}
$a^b$ by squaring: write $b$'s bits, square per position, multiply
into the answer on set bits --- $O(\log b)$ for exponents that would
take centuries linearly.
\begin{figure}[ht]\centering
% fig67-binpow.tex -- Figure 13.x: squaring tree for a^13; the exponent halves each level.
% fstyle.tex -- shared TikZ styles for the figures of this book (palette from lifeos.sty).
\tikzset{%
  cell/.style={draw=ink, minimum width=7mm, minimum height=6.5mm, inner sep=0pt, font=\footnotesize, fill=white},
  cellh/.style={cell, fill=accentsoft, draw=accent},
  cella/.style={cell, fill=amberbg, draw=amber},
  cellg/.style={cell, fill=goodbg, draw=good},
  cellx/.style={cell, fill=graybg, draw=slate, text=slate},
  vtx/.style={circle, draw=ink, minimum size=6.5mm, inner sep=0pt, font=\footnotesize, fill=white},
  vtxh/.style={vtx, draw=accent, fill=accentsoft},
  vtxa/.style={vtx, draw=amber, fill=amberbg},
  vtxg/.style={vtx, draw=good, fill=goodbg},
  vtxd/.style={vtx, draw=slate, fill=graybg, text=slate},
  lab/.style={font=\footnotesize, text=slate},
  labA/.style={font=\footnotesize, text=accent},
  labM/.style={font=\footnotesize, text=amber},
  idx/.style={font=\tiny, text=slate},
  arr/.style={-{stealth}, thick, ink},
  arrA/.style={-{stealth}, thick, accent},
  arrM/.style={-{stealth}, thick, amber},
  arrG/.style={-{stealth}, thick, good},
  edge/.style={thick, ink},
  edgeA/.style={thick, accent},
  edgeM/.style={thick, amber},
}%
% \cellrow{prefix}{x0}{y}{style}{v0,v1,...}: a row of cells named prefix0, prefix1, ... (pitch 0.75cm)
\providecommand{\cellrow}[5]{\foreach \v [count=\i from 0] in {#5} {\node[#4] (#1\i) at (#2+\i*0.75,#3) {\v};}}%
% \cellidx{prefix}{n}{start}: index labels under cells prefix0 .. prefix(n-1)
\providecommand{\cellidx}[3]{\foreach \i in {0,...,\numexpr#2-1\relax} {\node[idx, below=0pt of #1\i] {\the\numexpr\i+#3\relax};}}%
% \fitfig{<tikzpicture>}: shrink the picture to \linewidth only when it would overflow.
\ifdefined\fitbox\else\newsavebox{\fitbox}\fi%
\providecommand{\fitfig}[1]{\sbox{\fitbox}{#1}\ifdim\wd\fitbox>\linewidth\resizebox{\linewidth}{!}{\usebox{\fitbox}}\else\usebox{\fitbox}\fi}%

\begin{tikzpicture}[font=\footnotesize]
\tikzset{pn/.style={draw=slate, rounded corners=1.5pt, minimum height=8mm, minimum width=3.0cm, inner sep=2pt, fill=white}}
\tikzset{po/.style={pn, fill=amber!20, draw=amber}, pe/.style={pn, fill=accentsoft, draw=accent}, pz/.style={pn, fill=goodbg, draw=good}}
\foreach \e/\ex/\how/\y/\sty [count=\i] in {13/{a^{13}}/{$=a\cdot(a^{6})^{2}$\quad 13 odd: multiply by $a$}/0/po,
                                      6/{a^{6}}/{$=(a^{3})^{2}$\quad 6 even: just square}/-1.15/pe,
                                      3/{a^{3}}/{$=a\cdot(a^{1})^{2}$\quad 3 odd: multiply by $a$}/-2.3/po,
                                      1/{a^{1}}/{$=a\cdot(a^{0})^{2}$\quad 1 odd: multiply by $a$}/-3.45/po,
                                      0/{a^{0}}/{$=1$\quad base case}/-4.6/pz} {
  \node[\sty] (p\i) at (0,\y) {$\ex$};
  \node[lab, anchor=west] at (1.8,\y) {\how};
  \node[idx] at (-2.2,\y) {exp $\e$};
}
\foreach \i/\j in {1/2,2/3,3/4,4/5} \draw[arr, shorten >=1pt] (p\i.south) -- (p\j.north);
\node[labA, anchor=west] at (6.8,-4.6) {$\lceil\log_2 13\rceil=4$ squarings};
\end{tikzpicture}

\caption{Squaring tree for $a^{13}$. Notice only $O(\log b)$
nodes exist --- the exponent halves each level.}
\label{fig:binpow}
\end{figure}

\section{The \texttt{modint} struct}\label{sec:modint}\index{modint}
The class you write once: operators that reduce internally so
client code reads like ordinary arithmetic. Two silent failures
live inside: $(a-b)$ without $+m$ goes negative; $a*b$ in
\texttt{int} overflows 32-bit before the reduction.
\lstinputlisting{code/modint.h}

\section{Counting}\label{sec:counting}\index{stars and bars}\index{inclusion--exclusion}\index{derangement}\index{Catalan numbers}
Factorials plus inverse factorials give $\binom{n}{k} \bmod p$ in
$O(1)$ after $O(n)$ precompute. Stars and bars counts weak
compositions; inclusion--exclusion counts with forbidden events;
derangements and Catalan each reduce to named formulas.
\begin{example}[CSES \emph{Binomial Coefficients}]\label{ex:bincoef}
$\binom{a}{b} \bmod 10^9{+}7$ for up to $10^5$ queries.
\par\smallskip
\textbf{Strategy.} Precompute $f[n]=n!$ and $fi[n]=(n!)^{-1}$ once;
each query is three multiplications.
\par\smallskip
\begin{solution}
\lstinputlisting{code/binomial_coefficients.cpp}
Real output (CSES sample):
\begin{lstlisting}[numbers=none]
10
846527861
31465
\end{lstlisting}
\end{solution}
\end{example}
\begin{example}[CSES \emph{Distributing Apples}]\label{ex:apples}
Distribute $m$ identical apples among $n$ children.
\par\smallskip
\textbf{Strategy.} Lay the apples in a row and insert $n{-}1$
walls: each string of $n{+}m{-}1$ symbols is a distribution ---
$\binom{n+m-1}{n-1}$.
\par\smallskip
\begin{solution}
\lstinputlisting{code/distributing_apples.cpp}
Real output (CSES sample):
\begin{lstlisting}[numbers=none]
20
\end{lstlisting}
\end{solution}
\end{example}
\begin{figure}[ht]\centering
% fig68-starsbars.tex -- Figure 13.x: stars and bars; m = 5 apples among n = 3 children -> 7 symbols, 2 walls.
% fstyle.tex -- shared TikZ styles for the figures of this book (palette from lifeos.sty).
\tikzset{%
  cell/.style={draw=ink, minimum width=7mm, minimum height=6.5mm, inner sep=0pt, font=\footnotesize, fill=white},
  cellh/.style={cell, fill=accentsoft, draw=accent},
  cella/.style={cell, fill=amberbg, draw=amber},
  cellg/.style={cell, fill=goodbg, draw=good},
  cellx/.style={cell, fill=graybg, draw=slate, text=slate},
  vtx/.style={circle, draw=ink, minimum size=6.5mm, inner sep=0pt, font=\footnotesize, fill=white},
  vtxh/.style={vtx, draw=accent, fill=accentsoft},
  vtxa/.style={vtx, draw=amber, fill=amberbg},
  vtxg/.style={vtx, draw=good, fill=goodbg},
  vtxd/.style={vtx, draw=slate, fill=graybg, text=slate},
  lab/.style={font=\footnotesize, text=slate},
  labA/.style={font=\footnotesize, text=accent},
  labM/.style={font=\footnotesize, text=amber},
  idx/.style={font=\tiny, text=slate},
  arr/.style={-{stealth}, thick, ink},
  arrA/.style={-{stealth}, thick, accent},
  arrM/.style={-{stealth}, thick, amber},
  arrG/.style={-{stealth}, thick, good},
  edge/.style={thick, ink},
  edgeA/.style={thick, accent},
  edgeM/.style={thick, amber},
}%
% \cellrow{prefix}{x0}{y}{style}{v0,v1,...}: a row of cells named prefix0, prefix1, ... (pitch 0.75cm)
\providecommand{\cellrow}[5]{\foreach \v [count=\i from 0] in {#5} {\node[#4] (#1\i) at (#2+\i*0.75,#3) {\v};}}%
% \cellidx{prefix}{n}{start}: index labels under cells prefix0 .. prefix(n-1)
\providecommand{\cellidx}[3]{\foreach \i in {0,...,\numexpr#2-1\relax} {\node[idx, below=0pt of #1\i] {\the\numexpr\i+#3\relax};}}%
% \fitfig{<tikzpicture>}: shrink the picture to \linewidth only when it would overflow.
\ifdefined\fitbox\else\newsavebox{\fitbox}\fi%
\providecommand{\fitfig}[1]{\sbox{\fitbox}{#1}\ifdim\wd\fitbox>\linewidth\resizebox{\linewidth}{!}{\usebox{\fitbox}}\else\usebox{\fitbox}\fi}%

\begin{tikzpicture}[font=\footnotesize]
\foreach \s [count=\i from 0] in {a,a,w,a,w,a,a} {
  \if\s w \node[cell, minimum width=3mm, minimum height=9mm, fill=amber!30, draw=amber] (s\i) at (\i*0.8,0) {$|$};
  \else \node[cellh, minimum width=7mm, minimum height=7mm, rounded corners=3.5pt] (s\i) at (\i*0.8,0) {$\bullet$};\fi}
\draw[decorate,decoration={brace,mirror,raise=4pt},accent] (s0.south west) -- (s1.south east) node[midway,below=8pt,labA]{child 1: $2$};
\draw[decorate,decoration={brace,mirror,raise=4pt},accent] (s3.south west) -- (s3.south east) node[midway,below=8pt,labA]{child 2: $1$};
\draw[decorate,decoration={brace,mirror,raise=4pt},accent] (s5.south west) -- (s6.south east) node[midway,below=8pt,labA]{child 3: $2$};
\node[lab, anchor=west, align=left] at (6.3,0.4) {$m=5$ apples ($\bullet$) and $n-1=2$ walls ($|$)};
\node[lab, anchor=west, align=left] at (6.3,-0.3) {choose which $2$ of the $7$ slots are walls:};
\node[labM, anchor=west] at (6.3,-1.0) {$\binom{n+m-1}{n-1}=\binom{7}{2}=21$};
\end{tikzpicture}

\caption{The walls/children mapping. Notice the model works because
child-$i$'s count is just the run between walls --- no arithmetic
needed at all.}
\label{fig:starsbars}
\end{figure}
\begin{example}[CSES \emph{Christmas Party}]\label{ex:christmas}
Derangements of $n$ guests.
\par\smallskip
\textbf{Strategy.} Inclusion--exclusion over the ``fixed'' events,
or the equivalent recurrence
$D_n=(n-1)(D_{n-1}+D_{n-2})$ --- guest 1 picks a seat, and whether
that seat's owner lands in seat 1 splits the cases.
\par\smallskip
\begin{solution}
\lstinputlisting{code/christmas_party.cpp}
Real output (CSES sample):
\begin{lstlisting}[numbers=none]
9
\end{lstlisting}
\end{solution}
\end{example}
\begin{figure}[ht]\centering
% fig69-venn.tex -- Figure 13.x: inclusion--exclusion for three events.
% fstyle.tex -- shared TikZ styles for the figures of this book (palette from lifeos.sty).
\tikzset{%
  cell/.style={draw=ink, minimum width=7mm, minimum height=6.5mm, inner sep=0pt, font=\footnotesize, fill=white},
  cellh/.style={cell, fill=accentsoft, draw=accent},
  cella/.style={cell, fill=amberbg, draw=amber},
  cellg/.style={cell, fill=goodbg, draw=good},
  cellx/.style={cell, fill=graybg, draw=slate, text=slate},
  vtx/.style={circle, draw=ink, minimum size=6.5mm, inner sep=0pt, font=\footnotesize, fill=white},
  vtxh/.style={vtx, draw=accent, fill=accentsoft},
  vtxa/.style={vtx, draw=amber, fill=amberbg},
  vtxg/.style={vtx, draw=good, fill=goodbg},
  vtxd/.style={vtx, draw=slate, fill=graybg, text=slate},
  lab/.style={font=\footnotesize, text=slate},
  labA/.style={font=\footnotesize, text=accent},
  labM/.style={font=\footnotesize, text=amber},
  idx/.style={font=\tiny, text=slate},
  arr/.style={-{stealth}, thick, ink},
  arrA/.style={-{stealth}, thick, accent},
  arrM/.style={-{stealth}, thick, amber},
  arrG/.style={-{stealth}, thick, good},
  edge/.style={thick, ink},
  edgeA/.style={thick, accent},
  edgeM/.style={thick, amber},
}%
% \cellrow{prefix}{x0}{y}{style}{v0,v1,...}: a row of cells named prefix0, prefix1, ... (pitch 0.75cm)
\providecommand{\cellrow}[5]{\foreach \v [count=\i from 0] in {#5} {\node[#4] (#1\i) at (#2+\i*0.75,#3) {\v};}}%
% \cellidx{prefix}{n}{start}: index labels under cells prefix0 .. prefix(n-1)
\providecommand{\cellidx}[3]{\foreach \i in {0,...,\numexpr#2-1\relax} {\node[idx, below=0pt of #1\i] {\the\numexpr\i+#3\relax};}}%
% \fitfig{<tikzpicture>}: shrink the picture to \linewidth only when it would overflow.
\ifdefined\fitbox\else\newsavebox{\fitbox}\fi%
\providecommand{\fitfig}[1]{\sbox{\fitbox}{#1}\ifdim\wd\fitbox>\linewidth\resizebox{\linewidth}{!}{\usebox{\fitbox}}\else\usebox{\fitbox}\fi}%

\begin{tikzpicture}[font=\footnotesize]
\begin{scope}
  \fill[accent!25] (0,0) circle (1.1); \fill[amber!25] (1.3,0) circle (1.1); \fill[good!25] (0.65,-1.1) circle (1.1);
  \draw[accent, thick] (0,0) circle (1.1); \draw[amber, thick] (1.3,0) circle (1.1); \draw[good, thick] (0.65,-1.1) circle (1.1);
\end{scope}
\node[labA] at (-0.55,0.6) {$A$}; \node[labM] at (1.85,0.6) {$B$}; \node[font=\footnotesize, text=good] at (0.65,-1.75) {$C$};
\node[font=\scriptsize] at (0.65,0.25) {$A{\cap}B$}; \node[font=\scriptsize] at (0.1,-0.8) {$A{\cap}C$}; \node[font=\scriptsize] at (1.2,-0.8) {$B{\cap}C$};
\node[font=\scriptsize\bfseries] at (0.65,-0.35) {$ABC$};
\node[lab, anchor=west, align=left] at (3.6,0.5) {$|A\cup B\cup C| \;=\; |A|+|B|+|C|$};
\node[lab, anchor=west, align=left] at (4.5,-0.1) {$-\,|A{\cap}B|-|A{\cap}C|-|B{\cap}C|$};
\node[lab, anchor=west, align=left] at (4.5,-0.7) {$+\,|A{\cap}B{\cap}C|$};
\node[labM, anchor=west, align=left, text width=6cm] at (3.6,-1.6) {singles count each overlap twice (or thrice), pairs subtract too much, the triple adds the centre back};
\end{tikzpicture}

\caption{Inclusion--exclusion for three ``fixed'' events. Notice
the alternating signs repair the double and triple counting in
stages.}
\label{fig:venn}
\end{figure}
\begin{figure}[ht]\centering
\begin{tikzpicture}[font=\tiny,x=0.25cm,y=0.25cm]
\draw[slate!40,step=1] (0,0) grid (4,3);
\draw[-{stealth},amber,thick] (0,0)--(1,0)--(2,0)--(2,1)--(3,1)--(3,2)--(4,2)--(4,3);
\draw[dashed,slate] (0,0)--(4,3);
\node[font=\tiny] at (2,-1.2) {valid = never above diagonal};
\end{tikzpicture}

\caption{The ballot/reflection picture for Catalan counting.
Notice the diagonal is the legality boundary --- crossing it
identifies the bad paths that get reflected out.}
\label{fig:ballot}
\end{figure}

\section{Linear recurrences}\label{sec:matpow}\index{matrix exponentiation}
A linear recurrence is a matrix: $v_{n+1}=Mv_n$. Then $v_n=M^n v_0$
and binpow-on-matrices gives the $n$-th term in $O(\log n)$ ---
Fibonacci is the worked pattern, the same shape carries any
constant-coefficient recurrence.
\begin{example}[CSES \emph{Fibonacci Numbers}]\label{ex:fib}
$F_n \bmod 10^9{+}7$ for $n$ up to $10^{18}$.
\par\smallskip
\textbf{Strategy.} $M=\begin{pmatrix}1&1\\1&0\end{pmatrix}$ maps
$(F_k,F_{k-1})$ to $(F_{k+1},F_k)$; $F_n$ is the off-diagonal of
$M^n$.
\par\smallskip
\begin{solution}
\lstinputlisting{code/fibonacci.cpp}
Real output (CSES sample):
\begin{lstlisting}[numbers=none]
55
\end{lstlisting}
\end{solution}
\end{example}
\begin{figure}[ht]\centering
% fig71-matpow.tex -- Figure 13.x: Fibonacci by matrix power; M^10 = M^8 M^2 gives F_10 = 55.
% fstyle.tex -- shared TikZ styles for the figures of this book (palette from lifeos.sty).
\tikzset{%
  cell/.style={draw=ink, minimum width=7mm, minimum height=6.5mm, inner sep=0pt, font=\footnotesize, fill=white},
  cellh/.style={cell, fill=accentsoft, draw=accent},
  cella/.style={cell, fill=amberbg, draw=amber},
  cellg/.style={cell, fill=goodbg, draw=good},
  cellx/.style={cell, fill=graybg, draw=slate, text=slate},
  vtx/.style={circle, draw=ink, minimum size=6.5mm, inner sep=0pt, font=\footnotesize, fill=white},
  vtxh/.style={vtx, draw=accent, fill=accentsoft},
  vtxa/.style={vtx, draw=amber, fill=amberbg},
  vtxg/.style={vtx, draw=good, fill=goodbg},
  vtxd/.style={vtx, draw=slate, fill=graybg, text=slate},
  lab/.style={font=\footnotesize, text=slate},
  labA/.style={font=\footnotesize, text=accent},
  labM/.style={font=\footnotesize, text=amber},
  idx/.style={font=\tiny, text=slate},
  arr/.style={-{stealth}, thick, ink},
  arrA/.style={-{stealth}, thick, accent},
  arrM/.style={-{stealth}, thick, amber},
  arrG/.style={-{stealth}, thick, good},
  edge/.style={thick, ink},
  edgeA/.style={thick, accent},
  edgeM/.style={thick, amber},
}%
% \cellrow{prefix}{x0}{y}{style}{v0,v1,...}: a row of cells named prefix0, prefix1, ... (pitch 0.75cm)
\providecommand{\cellrow}[5]{\foreach \v [count=\i from 0] in {#5} {\node[#4] (#1\i) at (#2+\i*0.75,#3) {\v};}}%
% \cellidx{prefix}{n}{start}: index labels under cells prefix0 .. prefix(n-1)
\providecommand{\cellidx}[3]{\foreach \i in {0,...,\numexpr#2-1\relax} {\node[idx, below=0pt of #1\i] {\the\numexpr\i+#3\relax};}}%
% \fitfig{<tikzpicture>}: shrink the picture to \linewidth only when it would overflow.
\ifdefined\fitbox\else\newsavebox{\fitbox}\fi%
\providecommand{\fitfig}[1]{\sbox{\fitbox}{#1}\ifdim\wd\fitbox>\linewidth\resizebox{\linewidth}{!}{\usebox{\fitbox}}\else\usebox{\fitbox}\fi}%

\begin{tikzpicture}[font=\footnotesize]
\tikzset{mb/.style={draw=slate, rounded corners=2pt, inner sep=4pt, fill=white}}
\node[mb] (m1) at (0,0) {$M=\begin{pmatrix}1&1\\1&0\end{pmatrix}$};
\node[mb] (m2) at (3.5,0) {$M^2=\begin{pmatrix}2&1\\1&1\end{pmatrix}$};
\node[mb] (m4) at (7.0,0) {$M^4=\begin{pmatrix}5&3\\3&2\end{pmatrix}$};
\node[mb] (m8) at (10.5,0) {$M^8=\begin{pmatrix}34&21\\21&13\end{pmatrix}$};
\foreach \a/\b in {m1/m2,m2/m4,m4/m8} \draw[arrA] (\a.east) -- (\b.west) node[midway, above, labA, font=\scriptsize] {square};
\node[mb, draw=amber, fill=amberbg] (m10) at (7.0,-2.0) {$M^{10}=M^8M^2=\begin{pmatrix}89&55\\55&34\end{pmatrix}$};
\draw[arrM] (m8.south) -- (m10.east); \draw[arrM] (m2.south) -- (m10.west);
\node[labM, anchor=west] at (9.7,-2.0) {$F_{10}=55$};
\node[lab, align=center] at (6.0,-3.1) {exponentiate the transition: $\log_2 10$ squarings and one multiply; the state vector rides along};
\end{tikzpicture}

\caption{Matrix power. Notice we exponentiate the \emph{transition}
--- the state vector rides along for free.}
\label{fig:matpow}
\end{figure}

\begin{exercise}[medium]\label{exr:ext}
Given $ax+by=\gcd(a,b)$, express Bézout coefficients for $(b, a\bmod b)$
in terms of those for $(a,b)$ --- derive the two-line update.
\end{exercise}
\begin{exercise}[easy]\label{exr:primeinv}
Why does $a^{-1}\bmod m$ need $\gcd(a,m)=1$? Show what breaks
for $m$ composite and $a\mid m$.
\end{exercise}
\begin{exercise}[medium]\label{exr:model}
A problem asks for ordered $m$-tuples of nonnegative integers
summing to $n$. Which counting model applies, and what would
``at least one'' do to it?
\end{exercise}
\begin{exercise}[medium]\label{exr:matdim}
For $v_{n+1}=pv_{n}+qv_{n-1}+r$, write the transition matrix and
argue its dimension must be $3\times3$ --- what rides in the third slot?
\end{exercise}

\section{Chapter summary}
\begin{table}[ht]\centering
{\small
\begin{tabularx}{\linewidth}{@{}lYYY@{}}
\toprule
\tablehead\thead{Need} & \thead{Tool} & \thead{Cost} & \thead{Trap}\\
\midrule
gcd, inverse & ext.\ Euclid & $\log$ & sign on $(x,y)$\\
factorize & SPF sieve & $\log n$ & sieve bound\\
$a^b$ huge & binpow & $\log b$ & Fermat for exp-exp\\
all mod work & modint & native & neg.\ values, overflow\\
$n \choose k$ & fact$+$inv & $O(1)$ q & $p$ prime\\
weak comps.\ & stars-bars & formula & distinct $\ne$ identical\\
forbidden set & incl.--excl. & $2^k$ terms & signs alternate\\
linear recur.\ & matpow & $\log n$ & carry constants\\
\bottomrule
\end{tabularx}}
\caption{The problem's constraint names the model; the model names
the formula.}
\label{tab:ntcomb}
\end{table}

\section{Common mistakes}
\index{common mistakes}
\begin{itemize}
\item \textbf{Reducing the wrong exponent.} $a^{b^c}$ needs the
exponent reduced mod $m{-}1$, not mod $m$ --- the classic Exp-II
WA.
\item \textbf{Negative residues.} $(a-b)\,\%\,m$ can be negative in
C++; always add $m$ back before the final reduction.
\item \textbf{Inverse without coprimality.} Division $\bmod\ m$
silently fails when $\gcd(a,m)>1$ --- check before Fermat.
\item \textbf{Distinct vs identical objects.} Stars-and-bars counts
\emph{identical} objects; applying it to distinct objects misses the
permutations.
\item \textbf{Matrix dimension too small.} A recurrence with a
constant term needs the constant carried as an extra slot ---
$M^n v_0$ can't manufacture it.
\end{itemize}

\section{Additional exercises}
\begin{exercise}[easy] CSES \emph{Exponentiation}. \end{exercise}
\begin{exercise}[easy] CSES \emph{Counting Divisors}. \end{exercise}
\begin{exercise}[easy] CSES \emph{Common Divisors}. \end{exercise}
\begin{exercise}[easy] CSES \emph{Throwing Dice}. \end{exercise}
\begin{exercise}[medium] CSES \emph{Exponentiation II}. \end{exercise}
\begin{exercise}[medium] CSES \emph{Sum of Divisors}. \end{exercise}
\begin{exercise}[medium] CSES \emph{Binomial Coefficients}. \end{exercise}
\begin{exercise}[medium] CSES \emph{Distributing Apples}. \end{exercise}
\begin{exercise}[medium] CSES \emph{Creating Strings II}. \end{exercise}
\begin{exercise}[medium] CSES \emph{Permutation Rounds}. \end{exercise}
\begin{exercise}[medium] CSES \emph{Fibonacci Numbers}. \end{exercise}
\begin{exercise}[hard] CSES \emph{Prime Multiples} --- inclusion--exclusion. \end{exercise}
\begin{exercise}[hard] CSES \emph{Counting Coprime Pairs} --- M\"obius. \end{exercise}
\begin{exercise}[hard] CSES \emph{Christmas Party}. \end{exercise}
\begin{exercise}[hard] CSES \emph{Bracket Sequences I} --- Catalan. \end{exercise}
\begin{exercise}[hard] CSES \emph{Divisor Analysis}. \end{exercise}
\begin{exercise}[hard] CSES \emph{Graph Paths I/II} --- matpow on graphs. \end{exercise}

\chapternotes

\section*{Further reading}
\begin{itemize}
\item A.~Laaksonen, \emph{Competitive Programmer's Handbook}
(2017--2024), \url{https://cses.fi/book/book.pdf} --- free PDF;
the number-theory chapters this chapter is pitched from.
\item J.~Sannemo, \emph{Principles of Algorithmic Problem Solving}
(2025), \url{https://usaco.guide/PAPS1.pdf} --- free PDF; deeper
proofs for sieve bounds and modular machinery.
\item cp-algorithms, \emph{Algebra} and \emph{Combinatorics}
sections, \url{https://cp-algorithms.com/} --- the reference for
SPF variants, M"obius, and every named counting formula.
\item The USACO Guide, \emph{modular}, \emph{divisibility}, and
\emph{combo} modules, \url{https://usaco.guide/} --- the problem
lists this chapter's exercises were drawn from.
\item S.~Halim, F.~Halim, S.~Effendy, \emph{Competitive
Programming 4} (2020), \url{https://cpbook.net/} --- the widest
taxonomy of contest number theory and counting problems.
\item MIT OpenCourseWare, \emph{6.042J Mathematics for Computer
Science} (Fall 2010) --- the undergraduate course behind the
Watch picks below: \url{https://ocw.mit.edu/} (course 6.042J).
\item E.~W.~Weisstein (ed.), \emph{MathWorld},
\url{https://mathworld.wolfram.com/} --- formula lookup:
derangements, Catalan numbers, stars and bars, all with references.
\end{itemize}

\begin{rulebox}[Watch instead]
\begin{itemize}
\item \textbf{MIT OpenCourseWare --- ``Lec 4: Number Theory I''
(6.042J, Fall 2010)}, Marten van Dijk, 1h20m,
\url{https://www.youtube.com/watch?v=NuY7szYSXSw} --- covers
\S14.1--14.3: Euclid, gcd invariants, modular arithmetic.
(Verified 2026-10-11.)
\item \textbf{MIT OpenCourseWare --- ``Lec 16: Counting Rules I''
(6.042J, Fall 2010)}, Marten van Dijk, 1h20m,
\url{https://www.youtube.com/watch?v=pNt5Ll6hGqo} --- covers
\S14.6: bijection, product and sum rules behind stars-and-bars.
(Verified 2026-10-11.)
\item \textbf{Khan Academy --- ``Journey into Cryptography''
playlist}, Brit Cruise, short-clip playlist (segments
$\sim$10 min each),
\url{https://www.youtube.com/playlist?list=PLSQl0a2vh4HA50QhFIirlEZRXG4yjcoGM}
--- covers \S14.3: modular arithmetic through Fermat's little
theorem, the friendliest first pass. (Verified 2026-10-11.)
\end{itemize}
\emph{Notation note:} OCW writes $a \mid b$ and $\gcd$ as the
book does; Khan uses ``mod'' infix notation --- same semantics.
\end{rulebox}
The Euclidean and extended algorithms, sieve + SPF factorization,
and binary exponentiation follow Laaksonen~\cite{laaksonen2024} and
Sannemo~\cite{sannemo2025}; modular inverse, stars-and-bars,
derangement and Catalan treatments track cp-algorithms and the
USACO Guide~\cite{cpalgorithms,usacoguide}; counting-model choice
and the problem-band split follow Halim~\cite{halimcp4}; archive
problems are official CSES~\cite{cses}.
